\documentclass[11pt]{article}

\usepackage[utf8]{inputenc}
\usepackage[T1]{fontenc}
\usepackage[spanish]{babel}
\usepackage{amsmath,amssymb,mathtools,bm,graphicx,xcolor}
\usepackage{svg}
\svgsetup{inkscapelatex=false}
\usepackage[a4paper,margin=2.5cm]{geometry}
\usepackage[
    colorlinks=true,
    linkcolor=red,
    citecolor=green,
    urlcolor=red
]{hyperref}

\spanishdecimal{.}

\setlength{\parindent}{0pt}
\setlength{\parskip}{0.45em}
\setlength{\abovedisplayskip}{6pt}
\setlength{\belowdisplayskip}{6pt}
\setlength{\abovedisplayshortskip}{4pt}
\setlength{\belowdisplayshortskip}{4pt}

% Formato de encabezado y problemas
\newcommand{\encabezado}[3]{%
{\Large\bfseries #1\par}
\vspace{2pt}
{\normalsize #2\hfill #3\par}
\vspace{6pt}\hrule\vspace{8pt}}

\newcommand{\problema}[2]{%
\vspace{0.55em}
\noindent\textbf{#1.}\enspace{\small #2}\par
\vspace{0.3em}}

\newcommand{\inciso}[1]{%
\vspace{0.35em}
\noindent\textbf{(#1)}\enspace}

% Comandos auxiliares
\newcommand{\dd}{\,\mathrm{d}}
\newcommand{\ii}{\mathrm{i}}
\newcommand{\e}{\mathrm{e}}
\newcommand{\sen}{\operatorname{sen}}
\newcommand{\grad}{\nabla}
\newcommand{\laplaciano}{\nabla^2}
\newcommand{\R}{\mathbb{R}}

\begin{document}

\encabezado{Potencial de un cascarón esférico con densidad superficial}{Física Matemática II}{J.R., G.R.}

\problema{1}{Encuentre expresiones para el potencial electrostático $\phi(r,\theta,\varphi)$ dentro y fuera de un cascarón esférico de radio $R$, con densidad superficial de carga $\sigma(\theta,\varphi)$. Las condiciones sobre $r=R$ son
\begin{equation}
\label{continuidad-cascaron}
\phi_i(R^-,\theta,\varphi)=\phi_e(R^+,\theta,\varphi),
\end{equation}
\begin{equation}
\label{salto-cascaron}
-\varepsilon_0\left[\frac{\partial\phi_e}{\partial r}(R^+,\theta,\varphi)
-\frac{\partial\phi_i}{\partial r}(R^-,\theta,\varphi)\right]=\sigma(\theta,\varphi).
\end{equation}
Luego considere $\sigma(\theta,\varphi)=\zeta\sen^2\theta\sen\varphi\cos\varphi$.}

\textbf{Solución:} En las regiones $r<R$ y $r>R$ no hay carga de volumen. Por lo tanto, el potencial satisface la ecuación de Laplace
\begin{equation}
\label{laplace-cascaron}
\nabla^2\phi=0.
\end{equation}
La solución interior debe ser finita en el origen, mientras que la solución exterior debe anularse cuando $r\to\infty$. Por esto escribimos
\begin{align}
\label{potencial-interior-cascaron}
\phi_i(r,\theta,\varphi)
&=\sum_{\ell=0}^{\infty}\sum_{m=-\ell}^{\ell}
a_{\ell m}r^\ell Y_\ell^m(\theta,\varphi),
\qquad r<R,\\
\label{potencial-exterior-cascaron}
\phi_e(r,\theta,\varphi)
&=\sum_{\ell=0}^{\infty}\sum_{m=-\ell}^{\ell}
b_{\ell m}r^{-(\ell+1)}Y_\ell^m(\theta,\varphi),
\qquad r>R.
\end{align}

Imponiendo la continuidad (\ref{continuidad-cascaron}), tenemos
\begin{equation*}
\sum_{\ell,m}a_{\ell m}R^\ell Y_\ell^m(\theta,\varphi)
=\sum_{\ell,m}b_{\ell m}R^{-(\ell+1)}Y_\ell^m(\theta,\varphi).
\end{equation*}
Comparando coeficientes de los armónicos esféricos, obtenemos
\begin{equation}
\label{relacion-a-b-cascaron}
a_{\ell m}R^\ell=b_{\ell m}R^{-(\ell+1)}
\quad\Rightarrow\quad
b_{\ell m}=a_{\ell m}R^{2\ell+1}.
\end{equation}

Calculemos ahora las derivadas radiales. A partir de (\ref{potencial-interior-cascaron}) y (\ref{potencial-exterior-cascaron}),
\begin{align*}
\left.\frac{\partial\phi_i}{\partial r}\right|_{r=R}
&=\sum_{\ell,m}\ell a_{\ell m}R^{\ell-1}Y_\ell^m,\\
\left.\frac{\partial\phi_e}{\partial r}\right|_{r=R}
&=\sum_{\ell,m}-(\ell+1)b_{\ell m}R^{-(\ell+2)}Y_\ell^m.
\end{align*}
Sustituyendo (\ref{relacion-a-b-cascaron}) en la derivada exterior,
\begin{equation*}
\left.\frac{\partial\phi_e}{\partial r}\right|_{r=R}
=\sum_{\ell,m}-(\ell+1)a_{\ell m}R^{\ell-1}Y_\ell^m.
\end{equation*}
Luego, sustituyendo las derivadas radiales en la condición de salto (\ref{salto-cascaron}), resulta
\begin{equation}
\label{salto-en-serie-cascaron}
\varepsilon_0\sum_{\ell,m}(2\ell+1)a_{\ell m}R^{\ell-1}Y_\ell^m(\theta,\varphi)=\sigma(\theta,\varphi).
\end{equation}
Expandimos ahora la densidad superficial como
\begin{equation}
\label{sigma-expansion-cascaron}
\sigma(\theta,\varphi)=\sum_{\ell,m}\sigma_{\ell m}Y_\ell^m(\theta,\varphi).
\end{equation}
Comparando (\ref{salto-en-serie-cascaron}) con (\ref{sigma-expansion-cascaron}), encontramos
\begin{equation}
\label{coeficientes-generales-cascaron}
a_{\ell m}=\frac{\sigma_{\ell m}}{\varepsilon_0(2\ell+1)R^{\ell-1}},
\qquad
b_{\ell m}=\frac{\sigma_{\ell m}R^{\ell+2}}{\varepsilon_0(2\ell+1)}.
\end{equation}

Consideremos ahora el caso particular
\begin{equation}
\label{sigma-particular-cascaron}
\sigma(\theta,\varphi)=\zeta\sen^2\theta\sen\varphi\cos\varphi.
\end{equation}
Usando la identidad trigonométrica $\sen\varphi\cos\varphi=\sen(2\varphi)/2$, tenemos
\begin{equation*}
\sigma(\theta,\varphi)=\frac{\zeta}{2}\sen^2\theta\sen(2\varphi).
\end{equation*}
Por otra parte, los armónicos esféricos necesarios son
\begin{equation}
\label{armonicos-2pm2-cascaron}
Y_2^{\pm2}(\theta,\varphi)=\frac{1}{4}\sqrt{\frac{15}{2\pi}}\sen^2\theta\,\e^{\pm2\ii\varphi}.
\end{equation}
Además,
\begin{equation*}
\sen(2\varphi)=\frac{\e^{2\ii\varphi}-\e^{-2\ii\varphi}}{2\ii}.
\end{equation*}
Sustituyendo esta relación en (\ref{sigma-particular-cascaron}), y usando (\ref{armonicos-2pm2-cascaron}), se encuentra
\begin{equation}
\label{sigma-en-armonicos-cascaron}
\sigma(\theta,\varphi)
=-\ii\zeta\sqrt{\frac{2\pi}{15}}\left(Y_2^2-Y_2^{-2}\right).
\end{equation}
Por lo tanto, los únicos coeficientes no nulos son
\begin{equation}
\label{coef-sigma-particular-cascaron}
\sigma_{2,2}=-\ii\zeta\sqrt{\frac{2\pi}{15}},
\qquad
\sigma_{2,-2}=\ii\zeta\sqrt{\frac{2\pi}{15}}.
\end{equation}

Como los únicos términos no nulos tienen $\ell=2$, desde (\ref{coeficientes-generales-cascaron}) se obtiene
\begin{equation*}
a_{2,m}=\frac{\sigma_{2,m}}{5\varepsilon_0R},
\qquad
b_{2,m}=\frac{\sigma_{2,m}R^4}{5\varepsilon_0}.
\end{equation*}
Sustituyendo estos coeficientes en (\ref{potencial-interior-cascaron}) y (\ref{potencial-exterior-cascaron}), y volviendo a la forma real de la densidad, finalmente encontramos
\begin{equation*}
\boxed{
\phi_i(r,\theta,\varphi)=\frac{\zeta r^2}{5\varepsilon_0R}\sen^2\theta\sen\varphi\cos\varphi},
\qquad r<R,
\end{equation*}
\begin{equation*}
\boxed{
\phi_e(r,\theta,\varphi)=\frac{\zeta R^4}{5\varepsilon_0r^3}\sen^2\theta\sen\varphi\cos\varphi},
\qquad r>R.
\end{equation*}


\section*{Visualización numérica}

Para los mapas sobre la esfera se usa la latitud $\lambda=\pi/2-\theta$ en el eje vertical y la longitud $\varphi$ en el eje horizontal. El color representa el valor de la densidad o del potencial. Las líneas negras, cuando aparecen, son nodos de la función graficada, es decir, curvas donde ésta cambia de signo.

\begin{figure}[h]
\centering
\includesvg[width=0.86\textwidth]{figures/fig_01_sigma_esfera}
\caption{Mapa angular firmado de la densidad superficial $\sigma(\theta,\varphi)=\zeta\sen^2\theta\sen\varphi\cos\varphi$. El color indica el signo y la magnitud relativa de la carga superficial. Las líneas negras son nodos donde $\sigma(\theta,\varphi)=0$.}
\end{figure}

\begin{figure}[h]
\centering
\includesvg[width=0.86\textwidth]{figures/fig_05_superficie_phi_interior}
\caption{Mapa de magnitud angular asociado a la densidad superficial. Aquí el color muestra $|\sigma(\theta,\varphi)|$ normalizado, por lo que no distingue signo. La figura permite ubicar las regiones de mayor intensidad de la distribución sobre la esfera.}
\end{figure}

\begin{figure}[h]
\centering
\includesvg[width=0.70\textwidth]{figures/fig_02_potencial_interior_xy}
\caption{Potencial interior en el plano $z=0$. En estas unidades, $\phi_i\propto xy$, por lo que los ejes coordenados son líneas nodales. El color representa el signo y la magnitud relativa del potencial dentro del cascarón.}
\end{figure}

\begin{figure}[h]
\centering
\includesvg[width=0.70\textwidth]{figures/fig_03_potencial_exterior_xy}
\caption{Potencial exterior en el plano $z=0$. La estructura angular es la misma que en el interior, pero la amplitud disminuye al alejarse del cascarón debido al factor radial exterior.}
\end{figure}

\begin{figure}[h]
\centering
\includesvg[width=0.70\textwidth]{figures/fig_04_lineas_campo_xy}
\caption{Líneas del campo eléctrico exterior en el plano $z=0$, obtenidas a partir de $\vec E=-\nabla\phi$. El color representa el potencial y las curvas indican la dirección local del campo eléctrico en el corte.}
\end{figure}

\end{document}
